% META: {"id":"1SPE-DERLOCAL-CO-033","chapitre":"1SPE-DERIVATION-LOCAL","type_objet":"corrige","exercice_id":"1SPE-DERLOCAL-EX-033","status":"approved","origin":"generated","status_history":["generated","audited","accepted","approved"],"release_acceptance":"RELEASE_OWNER_BATCH_ACCEPTANCE","acceptance_closure_digest":"sha256:9b3ccf9a81c5520fbb7e03b7b2d3e7bf2057a02834b6d908bf3be5f49f3b553f","acceptance_receipt_digest":"sha256:4fad86f0282e0fa4e0419cca1ad6379ae7af5a280c04a4a90880f90a1c044242"}
% BEGIN-VERIFY
% from sympy import symbols, limit, expand
% h = symbols('h')
% taux = ((1 + h)**3 - 1) / h
% assert expand(taux) == 3 + 3*h + h**2
% assert limit(taux, h, 0) == 3
% END-VERIFY
\begin{corrige}{1SPE-DERLOCAL-EX-033}
On calcule $f(1+h)$ :
\[
f(1+h) = (1+h)^3 = 1 + 3h + 3h^2 + h^3
\]

On calcule $f(1)$ :
\[
f(1) = 1^3 = 1
\]

Le taux de variation est :
\[
\frac{f(1+h) - f(1)}{h} = \frac{1 + 3h + 3h^2 + h^3 - 1}{h} = \frac{3h + 3h^2 + h^3}{h}
\]

On factorise par $h$ au numérateur :
\[
\frac{h(3 + 3h + h^2)}{h} = 3 + 3h + h^2
\]

Lorsque $h$ tend vers $0$, l'expression $3 + 3h + h^2$ tend vers $3$.

Donc $\boxed{f'(1) = 3}$.
\end{corrige}
